% ============================================================
%  Teorema de Pitágoras — Apostila Premium | PratiqueJá
%  Engine: XeLaTeX
% ============================================================
\documentclass[12pt]{article}
\usepackage{../../pratiqueja}
\usepackage{tikz}
\usetikzlibrary{calc}

% \hlm — destaque de resultado em modo-math (cor sem bold)
\newcommand{\hlm}[1]{{\color{darkblue}#1}}

\tikzset{
  tri/.style={draw=darkblue, line width=1.2pt, fill=accentblue!12!white},
  trj/.style={draw=darkblue, line width=1.2pt, fill=badgegreen!12!white},
  hl/.style={draw=badgepurple!80!black, line width=1pt},
  ra/.style={draw=darkblue, line width=0.8pt},
  cl/.style={font=\footnotesize\bfseries},
}

\setsubject{Teorema de Pitágoras}

\begin{document}
\pagenumbering{arabic}
\setstretch{1.2}
\color{bodytext}

\heroheader{Teorema de Pitágoras}%
           {Catetos, Hipotenusa e Relações Métricas}%
           {Matemática · Intermediário}

\setlength{\columnsep}{18pt}
\setlength{\columnseprule}{0.5pt}
\def\columnseprulecolor{\color{exborder}}

\begin{multicols}{2}

%─────────────────────────────────────────────────────────────
\TW{Def}{badgepurple}{Definição}{Relação entre hipotenusa e catetos}

\regra{Num \textbf{triângulo retângulo}, o quadrado da
\textbf{hipotenusa} é igual à soma dos quadrados dos
\textbf{catetos}.}

\begin{center}\vspace{2pt}
\begin{tikzpicture}[scale=0.8]
\coordinate (A) at (0,0); \coordinate (B) at (3.6,0); \coordinate (C) at (0,2.6);
\draw[tri] (A)--(B)--(C)--cycle;
\draw[ra] (0.42,0) -- (0.42,0.42) -- (0,0.42);   % ângulo reto
\node[cl] at (1.8,-0.32){$b$};
\node[cl] at (-0.32,1.3){$c$};
\node[cl] at (2.1,1.55){$a$};
\node[font=\scriptsize] at (3.85,-0.05){};
\end{tikzpicture}
\end{center}

\formula{$a^2 = b^2 + c^2$}

\obs{$a$ = \textbf{hipotenusa} (oposta ao ângulo reto); $b, c$ =
\textbf{catetos} (formam o ângulo reto).}

%─────────────────────────────────────────────────────────────
\TW{Met}{darkblue}{Como Aplicar}{Isolar o lado desconhecido}

\regra{Hipotenusa desconhecida: $a = \sqrt{b^2 + c^2}$.\\[2pt]
Cateto desconhecido: $b = \sqrt{a^2 - c^2}$.}

\SH{Encontrando a hipotenusa}
\exemp{%
  $b=4,\, c=3$: \ $a^2 = 16 + 9 = 25 \Rightarrow a = \hlm{5}$\\[3pt]
  $b=5,\, c=12$: \ $a^2 = 25 + 144 = 169 \Rightarrow a = \hlm{13}$%
}

\SH{Encontrando um cateto}
\exemp{%
  $a=15,\, c=9$: \ $b^2 = 225 - 81 = 144 \Rightarrow b = \hlm{12}$\\[3pt]
  $a=10,\, c=6$: \ $b^2 = 100 - 36 = 64 \Rightarrow b = \hlm{8}$%
}

%─────────────────────────────────────────────────────────────
\TW{3,4,5}{badgegreen}{Ternas Pitagóricas}{Inteiros que satisfazem o teorema}

\regra{Uma \textbf{terna pitagórica} é um trio de inteiros
$(a,b,c)$ com $a^2 = b^2 + c^2$. Reconhecê-las evita calcular raízes.}

\exemp{%
  $(3,4,5)$ — $25 = 16 + 9$\\[2pt]
  $(5,12,13)$ — $169 = 144 + 25$\\[2pt]
  $(8,15,17)$ — $289 = 225 + 64$\\[2pt]
  $(7,24,25)$ — $625 = 576 + 49$%
}

\nota{Múltiplos também são ternas: $(6,8,10)$, $(9,12,15)$,
$(10,24,26)\ldots$}

%─────────────────────────────────────────────────────────────
\TW{$\Leftrightarrow$}{darkblue}{Recíproca e Classificação}{Identificar o tipo pelo maior lado}

\regra{Com lados $a \geq b \geq c$, compare $a^2$ com $b^2 + c^2$:}

\exemp{%
  $a^2 = b^2 + c^2$ → \hlm{retângulo}\\
  {\footnotesize $5^2 = 3^2 + 4^2 \Rightarrow 25 = 25$ \ok}\\[3pt]
  $a^2 > b^2 + c^2$ → \hlm{obtusângulo}\\
  {\footnotesize lados $4,5,7$: $49 > 41$ \ok}\\[3pt]
  $a^2 < b^2 + c^2$ → \hlm{acutângulo}\\
  {\footnotesize lados $4,5,6$: $36 < 41$ \ok}%
}

%─────────────────────────────────────────────────────────────
\TW{45°}{badgepurple}{Triângulos Notáveis}{45°-45°-90° e 30°-60°-90°}

\SH{Isósceles retângulo — 45°-45°-90°}
\begin{center}\vspace{1pt}
\begin{tikzpicture}[scale=0.8]
\coordinate (A) at (0,0); \coordinate (B) at (2.4,0); \coordinate (C) at (0,2.4);
\draw[trj] (A)--(B)--(C)--cycle;
\draw[ra] (0.4,0)--(0.4,0.4)--(0,0.4);
\node[cl] at (1.2,-0.32){$\ell$};
\node[cl] at (-0.32,1.2){$\ell$};
\node[cl] at (1.6,1.55){$\ell\sqrt{2}$};
\draw[graycolor, line width=0.7pt] ([shift=(135:0.5)]B) arc (135:180:0.5);
\node[font=\scriptsize] at ($(B)+(157.5:0.82)$){45°};
\draw[graycolor, line width=0.7pt] ([shift=(270:0.5)]C) arc (270:315:0.5);
\node[font=\scriptsize] at ($(C)+(292.5:0.72)$){45°};
\end{tikzpicture}
\end{center}
\exemp{%
  $a^2 = \ell^2 + \ell^2 = 2\ell^2 \;\Rightarrow\; a = \hlm{\ell\sqrt{2}}$\\
  {\footnotesize cateto $5$ → hipotenusa $5\sqrt{2} \approx 7{,}07$}%
}

\SH{Semi-equilátero — 30°-60°-90°}
\begin{center}\vspace{1pt}
\begin{tikzpicture}[scale=0.8]
\coordinate (A) at (0,0); \coordinate (B) at (3.0,0); \coordinate (C) at (0,1.732);
\draw[trj] (A)--(B)--(C)--cycle;
\draw[ra] (0.4,0)--(0.4,0.4)--(0,0.4);
\node[cl] at (1.5,-0.32){$\ell\sqrt{3}$};
\node[cl] at (-0.32,0.87){$\ell$};
\node[cl] at (1.8,1.05){$2\ell$};
\draw[graycolor, line width=0.7pt] ([shift=(150:0.6)]B) arc (150:180:0.6);
\node[font=\scriptsize] at ($(B)+(165:0.92)$){30°};
\draw[graycolor, line width=0.7pt] ([shift=(270:0.5)]C) arc (270:330:0.5);
\node[font=\scriptsize] at ($(C)+(300:0.7)$){60°};
\end{tikzpicture}
\end{center}
\exemp{%
  Relação $\;\ell : \ell\sqrt{3} : 2\ell\;$\\
  {\footnotesize cateto menor $3$ → cateto maior $3\sqrt{3}$, hipotenusa $6$}%
}

%─────────────────────────────────────────────────────────────
\TW{Rel}{darkblue}{Relações Métricas}{Altura sobre a hipotenusa}

\regra{A \textbf{altura} $h$ relativa à hipotenusa divide-a em duas
projeções $m$ e $n$ $(m+n=a)$. Daí decorrem quatro relações.}

\begin{center}\vspace{2pt}
\begin{tikzpicture}[scale=0.42]
\coordinate (A) at (0,0); \coordinate (B) at (10,0);
\coordinate (C) at (3.6,4.8); \coordinate (H) at (3.6,0);
\draw[tri] (A)--(B)--(C)--cycle;
\draw[hl] (C)--(H);
% ângulo reto em C
\draw[ra] (3.0,4.0)--(3.8,3.4)--(4.4,4.2);
% ângulo reto em H
\draw[ra] (3.6,0.7)--(4.3,0.7)--(4.3,0);
\node[cl] at (1.4,2.7){$c$};
\node[cl] at (7.2,2.9){$b$};
\node[cl] at (4.5,2.5){$h$};
\node[cl] at (1.8,-0.6){$m$};
\node[cl] at (6.8,-0.6){$n$};
\node[font=\scriptsize] at (3.6,-1.0){\scriptsize $H$};
\end{tikzpicture}
\end{center}

\formula{$h^2 = m\cdot n \quad\;\; b^2 = a\cdot m \quad\;\; c^2 = a\cdot n \quad\;\; a\cdot h = b\cdot c$}

\exemp{%
  \textbf{$b=6,\, c=8$ — achar $h, m, n$:}\\
  $a = \sqrt{36+64} = 10$\\
  $c^2 = a\cdot m \Rightarrow m = \dfrac{36}{10} = 3{,}6$\\
  $b^2 = a\cdot n \Rightarrow n = \dfrac{64}{10} = 6{,}4$\\
  $h^2 = m\cdot n = 23{,}04 \Rightarrow h = \hlm{4{,}8}$\\
  {\footnotesize Verif.: $a\cdot h = 48 = b\cdot c$ \ok}%
}

%─────────────────────────────────────────────────────────────
\TW{d}{badgegreen}{Distância entre Dois Pontos}{Aplicação no plano cartesiano}

\regra{A distância entre $(x_1,y_1)$ e $(x_2,y_2)$ vem de Pitágoras
aplicado às diferenças de coordenadas:}
\formula{$d = \sqrt{(x_2 - x_1)^2 + (y_2 - y_1)^2}$}

\exemp{%
  Entre $(15,14)$ e $(31,26)$:\\
  $d = \sqrt{16^2 + 12^2} = \sqrt{400} = \hlm{20}$\\[3pt]
  Entre $(1,1)$ e $(4,5)$:\\
  $d = \sqrt{9 + 16} = \sqrt{25} = \hlm{5}$%
}

%─────────────────────────────────────────────────────────────
\needspace{16\baselineskip}
\TW{Prob}{badgegreen}{Problema Contextualizado}{Triângulo retângulo oculto}

\exemp{%
  \textbf{P1 — escada:} escada de $5$\,m, pé a $3$\,m da parede.\\
  $5^2 = h^2 + 3^2 \Rightarrow h^2 = 16$\\
  Altura: $h = \hlm{4\text{ m}}$%
}

\columnbreak

\exemp{%
  \textbf{P2 — diagonal:} retângulo $6 \times 8$\,cm.\\
  $d^2 = 6^2 + 8^2 = 100$\\
  Diagonal: $d = \hlm{10\text{ cm}}$%
}

\end{multicols}

\end{document}
